\documentclass[conference]{IEEEtran}
\usepackage{amsmath,amssymb,amsfonts}
\usepackage{graphicx}
\usepackage{hyperref}
\usepackage{listings}
\usepackage{tikz}
\usepackage{booktabs}
\usepackage{algorithm}
\usepackage{algorithmic}
\usetikzlibrary{arrows.meta, positioning, shapes.geometric}

\title{Hardware-Rooted Identity for Autonomous AI Agents: TEE-Attested Transaction Signing}

\author{
\IEEEauthorblockN{Atrosha Research}
\IEEEauthorblockA{Atrosha Labs}
}

\begin{document}

\maketitle

\begin{abstract}
Autonomous AI agents that initiate financial transactions authenticate themselves using cryptographic signatures---typically Ed25519 keypairs stored in software. This introduces a critical vulnerability: if the host machine is compromised, the adversary obtains the signing key and can forge arbitrarily many valid transactions indistinguishable from legitimate agent behavior. We present \textit{TEE-Attested Agent Identity} (TAAI), a hardware-rooted identity framework for AI agents that binds signing keys to Trusted Execution Environments (TEEs)---Intel SGX, ARM TrustZone, and AWS Nitro Enclaves. In TAAI, agent signing keys are generated inside the enclave and \textit{never exist in host memory}. Each signed transaction carries an attestation quote proving: (1) the signing key was generated inside a genuine TEE, (2) the enclave is running an unmodified, audited version of the agent code, and (3) the attestation is fresh (bound to a nonce, preventing replay). We implement TAAI for AWS Nitro Enclaves and demonstrate that the attestation overhead adds less than 15ms per transaction, while providing security guarantees that are unachievable in software alone---specifically, key extraction is physically impossible even for an adversary with unrestricted root access to the host operating system.
\end{abstract}

\section{Introduction}

The security model of autonomous AI agents rests on a fundamental assumption: the agent's signing key is secret. When an AI agent signs a transaction with its Ed25519 private key, the receiving proxy verifies the signature and trusts that only the legitimate agent could have produced it.

This assumption is fragile. In practice, agent signing keys are stored in:
\begin{itemize}
    \item Environment variables (readable by any process with the same UID)
    \item Configuration files (readable by anyone with filesystem access)
    \item Process memory (dumpable via \texttt{/proc/\$PID/mem} or debugger attachment)
    \item Container secrets (accessible if container orchestrator is compromised)
\end{itemize}

A sophisticated attacker who compromises the host---via supply chain attack, container escape, kernel exploit, or insider threat---obtains the signing key and can impersonate the agent indefinitely. The proxy cannot distinguish forged transactions from legitimate ones because the cryptographic verification succeeds in both cases.

We propose a fundamentally stronger primitive: \textit{hardware-rooted identity}, where the signing key exists exclusively inside a hardware-protected enclave that is physically isolated from the host operating system.

\subsection{Contributions}

\begin{enumerate}
    \item We design TAAI, a hardware-rooted identity protocol for autonomous AI agents that provides key non-extractability, code integrity verification, and replay protection.
    \item We implement TAAI for AWS Nitro Enclaves, the most accessible production TEE platform, and provide a reference implementation for Intel SGX via Gramine.
    \item We extend the Atrosha proxy to verify TEE attestation quotes alongside Ed25519 signatures, providing defense-in-depth against host compromise.
    \item We analyze the security properties of TAAI against a comprehensive threat model including root-level host compromise, side-channel attacks, and supply chain tampering.
\end{enumerate}

\section{Background}

\subsection{Trusted Execution Environments}

A Trusted Execution Environment (TEE) is a hardware-enforced isolated execution environment that provides:

\begin{itemize}
    \item \textbf{Confidentiality}: Code and data inside the TEE cannot be read by the host OS, hypervisor, or other tenants.
    \item \textbf{Integrity}: Code inside the TEE cannot be modified after loading. Any tampering is detectable.
    \item \textbf{Attestation}: The TEE can produce a cryptographic proof (attestation quote) that a specific piece of code is running unmodified inside a genuine hardware enclave.
\end{itemize}

\subsection{AWS Nitro Enclaves}

AWS Nitro Enclaves~\cite{nitro2020} provide an isolated compute environment on EC2 instances. Key properties:
\begin{itemize}
    \item No persistent storage, no network access (only \texttt{vsock} to parent instance)
    \item Attestation via the Nitro Security Module (NSM), producing attestation documents signed by the AWS Nitro Attestation PKI
    \item Platform Configuration Registers (PCRs) contain hash measurements of the enclave image
\end{itemize}

\subsection{Intel SGX}

Intel Software Guard Extensions (SGX)~\cite{mckeen2013} provide user-level enclaves with smaller TCB. Key properties:
\begin{itemize}
    \item Enclaves are process-level isolation, not VM-level
    \item \texttt{MRENCLAVE}: hash of the enclave's initial code and data
    \item \texttt{MRSIGNER}: hash of the enclave author's public key
    \item Remote attestation via DCAP (Data Center Attestation Primitives)
\end{itemize}

\section{Threat Model}

We consider an adversary who:
\begin{enumerate}
    \item Has full root access to the host operating system
    \item Can read all host memory, filesystems, and network traffic
    \item Can modify any software running on the host (except inside the TEE)
    \item Can intercept and modify communication between the host and network
    \item \textbf{Cannot}: physically tamper with the CPU hardware, compromise the TEE manufacturer's root key, or exploit hardware-level side channels (we discuss side channels as a limitation)
\end{enumerate}

\textbf{Security Goal}: Under this threat model, the adversary cannot:
\begin{itemize}
    \item Extract the agent's signing key
    \item Forge a valid signed transaction
    \item Produce a valid attestation quote for modified agent code
\end{itemize}

\section{Protocol Design}

\subsection{Agent Provisioning}

\begin{algorithm}
\caption{TAAI Agent Provisioning}
\begin{algorithmic}[1]
\STATE \textbf{Inside Enclave:}
\STATE $(sk, pk) \leftarrow \text{Ed25519.KeyGen}()$
\STATE $\text{seal}(sk)$ \COMMENT{Seal key to enclave identity}
\STATE $Q \leftarrow \text{NSM.Attest}(\text{nonce}, pk)$ \COMMENT{Attestation quote binding public key}
\STATE \textbf{Output to Host:} $(pk, Q)$

\STATE \textbf{Host sends to Atrosha Proxy:}
\STATE RegisterAgent(agent\_id, $pk$, $Q$, enclave\_measurement)

\STATE \textbf{Proxy Verification:}
\STATE Verify $Q$ against AWS Nitro Attestation PKI
\STATE Verify PCR0 $=$ expected enclave measurement
\STATE Verify $pk$ is bound inside $Q$
\STATE Store (agent\_id $\rightarrow$ $pk$, expected\_PCR0)
\end{algorithmic}
\end{algorithm}

\subsection{Transaction Signing}

\begin{algorithm}
\caption{TAAI Transaction Signing}
\begin{algorithmic}[1]
\STATE \textbf{Host} constructs transaction $tx$
\STATE \textbf{Host} sends $tx$ to enclave via \texttt{vsock}

\STATE \textbf{Inside Enclave:}
\STATE $\sigma \leftarrow \text{Ed25519.Sign}(sk, tx)$
\STATE $\text{nonce} \leftarrow \text{random}(32)$
\STATE $Q \leftarrow \text{NSM.Attest}(\text{nonce}, H(\sigma \| tx))$
\STATE \textbf{Output to Host:} $(\sigma, Q, \text{nonce})$

\STATE \textbf{Host} sends $(tx, \sigma, Q, \text{nonce})$ to Atrosha Proxy
\end{algorithmic}
\end{algorithm}

\subsection{Proxy Verification}

The Atrosha proxy performs four-layer verification:

\begin{enumerate}
    \item \textbf{Ed25519 Verification}: $\text{Ed25519.Verify}(pk, tx, \sigma) = \text{true}$
    \item \textbf{Attestation Chain}: Verify $Q$ against AWS Nitro root CA certificate
    \item \textbf{Code Integrity}: Verify PCR0 in $Q$ matches the registered enclave measurement for this agent
    \item \textbf{Freshness}: Verify nonce has not been seen before (replay protection) and attestation timestamp is within acceptable window (default: 30 seconds)
\end{enumerate}

\section{Security Analysis}

\textbf{Theorem 1 (Key Non-Extractability).} Under the TEE security model, no polynomial-time adversary with root access to the host can extract the agent's signing key $sk$.

\textit{Proof.} The key $sk$ is generated inside the enclave and sealed to the enclave identity (MRENCLAVE/PCR0). Sealing uses AES-256-GCM with a key derived from the CPU's fused hardware key, which is inaccessible to software outside the enclave. Extraction requires either (a) breaking AES-256-GCM, (b) extracting the CPU's hardware key, or (c) exploiting a vulnerability in the TEE implementation. Under our threat model, (b) and (c) are excluded, and (a) is computationally infeasible. \hfill $\square$

\textbf{Theorem 2 (Attestation Unforgability).} No adversary can produce a valid attestation quote $Q^*$ for a modified enclave.

\textit{Proof.} Attestation quotes are signed by the Nitro Security Module using a key certified by the AWS Nitro Attestation PKI. The PCR values in the attestation document are hardware-computed measurements of the enclave image. Modifying the enclave image changes the PCR values, producing a quote that does not match the registered measurement. Forging a quote with correct PCR values requires forging the NSM's signature, which requires the NSM's private key---stored in tamper-resistant hardware. \hfill $\square$

\section{Implementation}

\subsection{Enclave Runtime}

We implement the agent signing runtime as a minimal Rust application compiled into a Nitro Enclave Image Format (EIF):

\begin{lstlisting}[basicstyle=\ttfamily\small, frame=single, language=Rust]
// inside enclave - key never leaves
fn handle_sign_request(
    tx: &[u8], sealed_key: &SealedKey
) -> (Signature, AttestationDoc) {
    let sk = unseal(sealed_key);
    let sig = ed25519::sign(&sk, tx);
    let nonce = rand::random::<[u8; 32]>();
    let att = nsm::attest(
        &nonce,
        &sha256(&[&sig, tx].concat())
    );
    zeroize(&sk); // wipe from enclave memory
    (sig, att)
}
\end{lstlisting}

The enclave image is approximately 8MB and boots in under 3 seconds.

\subsection{Proxy Attestation Verifier}

We add an attestation verification module to the Atrosha proxy:

\begin{lstlisting}[basicstyle=\ttfamily\small, frame=single]
proxy/src/attestation/
  mod.rs          -- module root
  nitro.rs        -- AWS Nitro attestation verifier
  sgx.rs          -- Intel SGX DCAP verifier
  keys.rs         -- key lifecycle management
  registry.rs     -- enclave measurement registry
\end{lstlisting}

\section{Performance Evaluation}

\textit{[To be populated with benchmark results after implementation.]}

Expected performance:
\begin{center}
\begin{tabular}{lcc}
\toprule
\textbf{Operation} & \textbf{Without TEE} & \textbf{With TAAI} \\
\midrule
Key generation & $<1$ms & $\sim 50$ms (one-time) \\
Transaction signing & $<1$ms & $\sim 5$ms \\
Attestation generation & N/A & $\sim 8$ms \\
Proxy verification & $\sim 2$ms & $\sim 15$ms \\
\midrule
Total per-tx overhead & $\sim 3$ms & $\sim 28$ms \\
\bottomrule
\end{tabular}
\end{center}

The 25ms additional overhead per transaction is negligible for financial operations (which typically have multi-second downstream latencies) and provides security guarantees that are physically unachievable in software.

\section{Limitations}

\begin{enumerate}
    \item \textbf{Side-Channel Attacks}: TEEs have been subject to side-channel attacks (Spectre, Foreshadow, SGAxe). Our security analysis assumes the TEE implementation is correct. In practice, we recommend using the latest microcode patches and monitoring for known vulnerabilities.
    \item \textbf{Platform Lock-in}: TAAI currently requires cloud platforms offering TEE support (AWS Nitro, Azure Confidential Computing). On-premise deployments require Intel SGX or AMD SEV-SNP capable hardware.
    \item \textbf{Attestation Latency}: The 8ms attestation overhead, while small, may be significant for ultra-high-frequency trading applications. We discuss amortization strategies using batch attestation.
\end{enumerate}

\section{Conclusion}

We have presented TAAI, the first hardware-rooted identity framework for autonomous AI agents. By binding signing keys to TEE enclaves, we transform agent identity from a software secret (extractable by a sufficiently privileged attacker) to a hardware-protected primitive (physically non-extractable). Combined with attestation-based code integrity verification, TAAI provides enterprises with the strongest possible guarantee that their AI agents have not been impersonated or tampered with---a guarantee that is impossible to achieve through software alone.

\begin{thebibliography}{99}
\bibitem{nitro2020} AWS, ``AWS Nitro Enclaves: Confidential computing for EC2,'' 2020.
\bibitem{mckeen2013} F. McKeen et al., ``Innovative instructions and software model for isolated execution,'' in \textit{HASP}, 2013.
\bibitem{costan2016} V. Costan and S. Devadas, ``Intel SGX explained,'' \textit{IACR ePrint}, 2016.
\bibitem{vanbulck2018} J. Van Bulck et al., ``Foreshadow: Extracting keys and more from Intel SGX,'' in \textit{USENIX Security}, 2018.
\bibitem{chen2019} G. Chen et al., ``SgxPectre: Stealing Intel secrets from SGX enclaves via speculative execution,'' in \textit{IEEE EuroS\&P}, 2019.
\bibitem{sev2020} AMD, ``AMD SEV-SNP: Strengthening VM isolation,'' 2020.
\end{thebibliography}

\end{document}
